\subsection{\problemblock{*Define} Data Block}\label{sec:trajectory-block-define}

The standard trajectory variables computed in TAOS are listed in the appendix.
These variables are always computed and are available for printout and plotting;
however, some problems require quantities that are not among the standard output
variables. The required quantities differ with the type of problem.

TAOS provides user-defined output variables for this purpose. As long as a variable
of interest can be calculated from the standard variables with one or more equations,
that is, as long as it is a function of the standard variables, it can be computed as a
user-defined variable. A user-defined variable is given a name that must be different
from every standard output-variable name. The name is followed by one or more
equations showing how the new variable is computed. The equations may compute
the variable directly, or they may define its time derivative so that the variable is
integrated along with the equations of motion.

The equations are entered in a simplified C language. The syntax is also similar to
Fortran, so users familiar with either language should have little difficulty entering
the equations. The following examples illustrate the input form. The first example,

\begin{lstlisting}[style=taosmanual]
*define alt-km
  alt-km = alt / 3280.84;
\end{lstlisting}

computes altitude in kilometers rather than feet. The name of the new variable,
\texttt{alt-km}, follows the \problemblock{*define} keyword, and it must appear on
the left side of an assignment somewhere in the definition. This example contains
only one simple equation, so \texttt{alt-km} is the variable on the left side of the
equal sign. Every value and variable on the right side must be a standard output
variable, a previously input user-defined variable, a temporary variable, or a
constant. Here \statevar{alt} is a standard output variable and 3280.84 is a constant.
Each statement ends with a required semicolon.

\taossubhead{Temporary Variables}

For a blunt body of revolution in hypersonic flow, the stagnation-point heat-transfer
rate can be approximated by

\begin{equation}
  q_s = \frac{17600}{\sqrt{R_n}}
        \sqrt{\frac{\rho}{\rho_0}}
        \left(\frac{\lVert\vec V\rVert}{\lVert\vec V_c\rVert}\right)^{3.15}
        \quad \mathrm{Btu/(ft^2\,s)},
  \label{eq:stagnation-heating}
\end{equation}

where $R_n$ is the nose radius, $\rho$ and $\rho_0$ are the ambient and sea-level
densities, $\lVert\vec V\rVert$ is flight velocity, and
$\lVert\vec V_c\rVert$ is satellite velocity, usually 26,000~ft/s. The following data
block defines a user variable, \texttt{stag-heat}, from this equation:

\begin{lstlisting}[style=taosmanual]
*define stag-heat
  rnose = 1.0 / 12.0;
  stag-heat = (17600.0/sqrt(rnose))*sqrt(rho/0.0023769)
              * (vel/26000.0)^3.15;
\end{lstlisting}

A one-inch nose radius is converted to feet and stored in the temporary variable
\texttt{rnose}. Temporary variables are used only while the user-defined variable is
being calculated. If a name is not a standard output variable or a previously input
user-defined variable, TAOS assumes it is a temporary variable. A temporary
variable must be assigned before it is used.

The only unusual operator in the stagnation-heating equation is
\texttt{\textasciicircum}, which raises a value to a power. TAOS uses this operator
rather than the C \texttt{pow} function or the Fortran \texttt{**} operator. The final
assignment occupies two input lines. A statement continues until a semicolon is
encountered; as in C, the semicolon is required at the end of every statement.

\taossubhead{Math Operations and Functions}

\Cref{tab:user-defined-math} lists the operators and functions that may appear in
user-defined-variable equations. The standard arithmetic and relational operators
have the same precedence they have in C and Fortran. The mathematical functions
are defined with the usual C-language meanings.

\begin{longtable}{@{}>{\ttfamily}lp{0.20\linewidth}>{\ttfamily}lp{0.37\linewidth}@{}}
  \caption{User-defined-variable math operations and functions.}
  \label{tab:user-defined-math}\\
  \toprule
  \normalfont Operator & Description & \normalfont Function & Description \\
  \midrule
  \endfirsthead
  \toprule
  \normalfont Operator & Description & \normalfont Function & Description \\
  \midrule
  \endhead
  +   & Add & sin(x)   & Sine of $x$ in degrees \\
  -   & Subtract & cos(x)   & Cosine of $x$ in degrees \\
  *   & Multiply & tan(x)   & Tangent of $x$ in degrees \\
  /   & Divide & asin(x)  & Inverse sine, $x\in[-1,1]$ \\
  \textasciicircum{} & Exponentiation & acos(x)  & Inverse cosine, $x\in[-1,1]$ \\
  >   & Greater than & atan(x)  & Inverse tangent in $[-\pi/2,\pi/2]$ \\
  <   & Less than & atan2(x) & Inverse tangent in $[-\pi,\pi]$ \\
  >=  & Greater than or equal & sinh(x)  & Hyperbolic sine \\
  <=  & Less than or equal & cosh(x)    & Hyperbolic cosine \\
  ==  & Equal to & tanh(x)    & Hyperbolic tangent \\
  !=  & Not equal to & exp(x)     & Exponential function $e^x$ \\
  \&\& & Logical and & log(x)     & Natural logarithm \\
  ||  & Logical or & log10(x)   & Base-10 logarithm \\
      & & sqrt(x)    & Square root \\
      & & ceil(x)    & Smallest integer not less than $x$ \\
      & & floor(x)   & Largest integer not greater than $x$ \\
      & & abs(x)     & Absolute value \\
      & & table(id)  & Lookup from an \tabletype{output} table named \texttt{id} \\
  \bottomrule
\end{longtable}

The \texttt{table} function is special because it retrieves a value from a table file.
Its argument is the table-identification name described in \cref{sec:simple-tables},
and the table must have type \tabletype{output}. For example,

\begin{lstlisting}[style=taosmanual]
*define special
  special = (rho/0.0023769)*table(special-tbl);
\end{lstlisting}

references the output table \texttt{special-tbl}. One possible use is aerodynamic
heating analysis: output tables containing heat-transfer rates can be interpolated
along a trajectory, and the resulting value can then be printed or limited by an
optimization constraint.

If-then statements control how a user-defined variable is calculated. For example,

\begin{lstlisting}[style=taosmanual]
*define nx-max
  nx-max = max(nx);
  if (nx-max < 0) nx-max = 0;
\end{lstlisting}

calculates maximum axial loading and limits \texttt{nx-max} to positive values. The
relationship in an \texttt{if} statement is enclosed in parentheses. When the
relationship is true, the following statement is executed.

More than one statement can be executed by enclosing the statements in curly
brackets:

\begin{lstlisting}[style=taosmanual]
*define new-var
  new-var = 0;
  if (time > 20) {
    ratio = pres / 2116.2;
    new-var = ratio * table(tbl-id);
  }
\end{lstlisting}

C uses braces to group statements; the Fortran keywords \texttt{then} and
\texttt{endif} are not used. Braces are also required for an if-then-else statement:

\begin{lstlisting}[style=taosmanual]
*define var
  if (mach > 1.0) {
    x = 1 + mach*mach;
    var = sqrt(x) * table(outx);
  } else {
    var = table(outy);
  }
\end{lstlisting}

The current version of TAOS supports assignment statements and if-then-else
statements only. While loops, for loops, and other C control statements cannot be
used.

\taossubhead{Integral Variables}

A common use of an integral user-defined variable is a flight-path optimization
constraint. Constraints at particular points in a trajectory are entered directly in the
\problemblock{*optimize} block, but an integral variable is needed when the entire
flight path must remain within a limit.

For example, a minimum dynamic pressure of 100~lbf/ft\textsuperscript{2} can be
represented by

\begin{align}
  q>100\ \text{or}\ v<1000 &: \qquad
    q_{\min}=\int 0\,dt, \label{eq:qmin-satisfied}\\
  \text{otherwise} &: \qquad
    q_{\min}=\int (100-q)^2\,dt. \label{eq:qmin-violated}
\end{align}

Optimization constraints work best when the constraint variable is less than zero
while the constraint is satisfied and greater than zero when it is violated. If dynamic
pressure is above the minimum while the vehicle is moving rapidly, the derivative is
zero. If dynamic pressure falls below 100~lbf/ft\textsuperscript{2}, the derivative is
positive. The definition is entered as

\begin{lstlisting}[style=taosmanual]
*define integral qmin=-0.10
  if (dynprs > 100 || vel < 1000) {
    qmin = 0.0;
  } else {
    qmin = (100.0 - dynprs)^2;
  }
\end{lstlisting}

The \texttt{integral} keyword says that \texttt{qmin} is a derivative to be integrated
with respect to time along the flight path. Its initial value is $-0.10$, so the
constraint is initially satisfied. The conditional statement means: if dynamic
pressure is greater than 100 or velocity is less than 1000, set the derivative to zero;
otherwise, set it equal to the square of 100 minus dynamic pressure.
